\documentclass{rapportFSR}
\usepackage[utf8]{inputenc}
\usepackage{lipsum}
\usepackage{minted}
\usepackage{mdframed}
\usepackage{tcolorbox}
\usepackage{pifont}
\usepackage{fontawesome5}
\usepackage{listings}
\usepackage{xcolor}
\usepackage{booktabs}   % Pour des lignes horizontales de qualité professionnelle
\usepackage{tabularx}   % Pour un tableau à largeur fixe
\usepackage{graphicx}   % Si vous souhaitez ajouter des icônes (optionnel)
\usepackage{colortbl}
\setlength{\parindent}{0pt} % Supprime l'alinéa pour tout le document

\usepackage{listings}
\usepackage{xcolor}
\usepackage{enumitem}
% Définition de couleurs modernes (Style Nord/VS Code)
\definecolor{pygreen}{RGB}{46, 139, 87}
\definecolor{pyblue}{RGB}{0, 112, 192}
\definecolor{pyred}{RGB}{200, 40, 40}
\definecolor{pypurple}{RGB}{128, 0, 128}
\definecolor{backcolour}{RGB}{248, 248, 248}
\definecolor{graynumbers}{RGB}{128, 128, 128}

\lstset{
    language=Python,
    backgroundcolor=\color{backcolour},   
    commentstyle=\color{pygreen}\itshape,
    keywordstyle=\color{pyblue}\bfseries,
    numberstyle=\tiny\color{graynumbers},
    stringstyle=\color{pyred},
    basicstyle=\ttfamily\small,
    %
    % Gestion des accents et caractères spéciaux (essentiel pour les rapports en français)
    extendedchars=true,
    literate={á}{{\'a}}1 {é}{{\'e}}1 {í}{{\'i}}1 {ó}{{\'o}}1 {ú}{{\'u}}1 {Á}{{\'A}}1 {É}{{\'E}}1 {Í}{{\'I}}1 {Ó}{{\'O}}1 {Ú}{{\'U}}1 {à}{{\`a}}1 {è}{{\`e}}1 {ì}{{\`i}}1 {ò}{{\`o}}1 {ù}{{\`u}}1 {À}{{\`A}}1 {È}{{\`E}}1 {Ì}{{\`I}}1 {Ò}{{\`O}}1 {Ù}{{\`U}}1 {ä}{{\"a}}1 {ë}{{\"e}}1 {ï}{{\"i}}1 {ö}{{\"o}}1 {ü}{{\"u}}1 {Ä}{{\"A}}1 {Ë}{{\"E}}1 {Ï}{{\"I}}1 {Ö}{{\"O}}1 {Ü}{{\"U}}1 {â}{{\^a}}1 {ê}{{\^e}}1 {î}{{\^i}}1 {ô}{{\^o}}1 {û}{{\^u}}1 {Â}{{\^A}}1 {Ê}{{\^E}}1 {Î}{{\^I}}1 {Ô}{{\^O}}1 {Û}{{\^U}}1 {œ}{{\oe}}1 {Œ}{{\OE}}1 {æ}{{\ae}}1 {Æ}{{\AE}}1 {ç}{{\c{c}}}1 {Ç}{{\c{C}}}1 {ø}{{\o}}1 {å}{{\r a}}1 {Å}{{\r A}}1 {ñ}{{\~n}}1 {Ñ}{{\~N}}1 {¿}{{\textquestiondown}}1 {¡}{{\textexclamdown}}1,
    %
    breaklines=true,
    breakatwhitespace=true,
    showstringspaces=false,
    showspaces=false,
    showtabs=false,
    numbers=left,                    
    numbersep=10pt,                  
    frame=lines, % Lignes en haut et en bas seulement pour un look plus moderne
    rulecolor=\color{black!20},
    captionpos=b, % Légende en bas
    keepspaces=true,                 
    tabsize=4
}

\definecolor{lightblue}{RGB}{220,235,250}


\newtcolorbox{questionbox}{
colback=lightblue,
colframe=blue,
boxrule=0.8pt,
arc=3mm,
left=6pt,
right=6pt,
top=6pt,
bottom=6pt
}



\definecolor{bg}{rgb}{0.95,0.95,0.95} 

\title{Rapport FSR - Template}

\begin{document}

%------------------ Informations du projet (à remplir) ------------------
\titre{Titre du projet sur une ou deux lignes} 
\sujet{Rapport de projet} 
\enseignant{Pr. Prénom \textsc{NOM}} 
\eleves{Prénom \textsc{NOM} \\Prénom \textsc{NOM}
} 


\fairemarges
\fairepagedegarde
\tabledematieres
\newpage
\listoffigures
\newpage
\listoftables
\newpage

\section*{Glossaire}

\begin{table}[h]
\centering
\begin{tabular}{|l|p{10cm}|}
\hline
\textbf{Terme} & \textbf{Définition} \\
\hline
Terme 1 & Définition du terme. \\
\hline
Terme 2 & Définition du terme. \\
\hline
Terme 3 & Définition du terme. \\
\hline
Terme 4 & Définition du terme. \\
\hline
\end{tabular}
\end{table}


 \begin{center}
     \href{https://github.com/}{\color{blue}Lien du projet en GitHub}
 \end{center}


\newpage



\section{Contexte Général et Problématique}
\subsection{Contexte du projet}
Présentez le domaine du projet et le cadre dans lequel il s’inscrit.

Le contexte technique repose sur trois piliers :

\begin{itemize}[leftmargin=*,label=$\bullet$]
    \item \textbf{Pilier 1 :} Description.
    \item \textbf{Pilier 2 :} Description.
    \item \textbf{Pilier 3 :} Description.
\end{itemize}

\subsection{Problématique}
Exposez les problèmes que le projet cherche à résoudre :
\begin{itemize}[leftmargin=*, label=$\bullet$]
    \item \textbf{Problème 1 :} Description.
    \item \textbf{Problème 2 :} Description.
    \item \textbf{Problème 3 :} Description.
\end{itemize}

\subsection{Objectifs du projet}
Décrivez l’objectif principal du projet et les objectifs secondaires.

% --- Chapitre 2 ---
\section{État de l'Art et Étude Théorique}
\subsection{Notion ou technologie 1}
Définition et rôle dans le projet \cite{ref1}.

\begin{figure}[H]
    \centering
    \includegraphics[width=0.3\linewidth]{capture.png}
    \caption{Titre de la figure}
    \label{fig:notion1}
\end{figure}

\begin{table}[h]
\centering
\begin{tabular}{ll}
\toprule
Élément & Rôle dans le projet \\
\midrule
Élément 1 & Rôle \\
Élément 2 & Rôle \\
Élément 3 & Rôle \\
\bottomrule
\end{tabular}
\caption{Titre du tableau}
\end{table}

\subsubsection{Sous-notion}

Explication détaillée.

\subsection{Notion ou technologie 2}
Définition et rôle dans le projet.

% --- Chapitre 3 ---
\section{Conception et Architecture Technique}
Présentez la phase de conception de la solution.

\subsection{Architecture de la solution}

\begin{figure}[H]
    \centering
    \includegraphics[width=1\linewidth]{architecture.png}
    \caption{Architecture de la solution}
    \label{fig:architecture}
\end{figure}

\subsection{Configuration}
L’extrait suivant présente la configuration utilisée :

\begin{lstlisting}[language=Python, caption={Extrait de code}]
# Votre code ici
def main():
    print("Rapport FSR")
\end{lstlisting}

Commentez l’extrait de code.

% --- Chapitre 4 ---
\section{Réalisation Pratique et Tests}
\subsection{Déploiement}

\begin{lstlisting}[language=bash, caption={Commande de déploiement}]
commande --option
\end{lstlisting}

\begin{figure}[H]
    \centering
    \includegraphics[width=0.75\linewidth]{capture.png}
    \caption{Titre de la capture}
    \label{fig:deploiement}
\end{figure}

\subsection{Tests et validation}
Décrivez les tests réalisés et leurs résultats.

% --- Chapitre 5 ---
\section{Analyse Comparative et Bilan}

\subsection{Comparaison des approches}

\begin{table}[h]
\centering
\begin{tabular}{|l|p{5cm}|p{5cm}|}
\hline
\textbf{Indicateur} & \textbf{Approche classique} & \textbf{Approche du projet} \\ \hline
Indicateur 1 & Valeur & Valeur \\ \hline
Indicateur 2 & Valeur & Valeur \\ \hline
Indicateur 3 & Valeur & Valeur \\ \hline
\end{tabular}
\caption{Comparaison des approches}
\end{table}

\subsection{Bilan Technique}

Résumez les apports du projet.

\subsection{Conclusion et Perspectives}

Concluez et présentez les perspectives.
\newpage


\begin{thebibliography}{9}

\bibitem{ref1}
Auteur, \emph{Titre du document}, année.  
\url{https://exemple.com}

\bibitem{ref2}
Auteur, \emph{Titre du document}, année.  
\url{https://exemple.com}

\end{thebibliography}

\end{document}
